\documentclass[10pt,a4paper]{amsart}
\usepackage[margin=19mm]{geometry}
\usepackage{amsmath,amssymb,mathtools}
\usepackage[scr=rsfs,cal=euler]{mathalfa}
\usepackage{xfrac}
\usepackage[unicode,hidelinks,bookmarks=false]{hyperref}
\newcommand{\ie}{i.e.\ }
\newcommand{\cf}{cf.\ }
\newcommand{\resp}{respectively\ }
\newcommand{\Subsection}{\subsection}
\providecommand{\nobreakdash}{\nobreak\hbox{-}}
\setlength{\emergencystretch}{2em}
\multlinegap0pt
\usepackage{amssymb}
\usepackage{mathtools}
\usepackage{braket}
\usepackage{bm}
\usepackage[shortlabels]{enumitem}
\usepackage{xcolor}
\newtheorem{assumption}{Assumption}
\newtheorem{proposition}{Proposition}
\def\RR{\mathbb{R}} %reals
\title{Martin boundary and invariant fields for multiplicative SHE}
\author{Hongyi Chen}
\date{Source: arXiv:2602.16126v5 (CC BY 4.0)}
\begin{document}
\maketitle

\noindent\emph{Source and licence.} Martin boundary and invariant fields for multiplicative SHE. Version arXiv:2602.16126v5 (CC BY 4.0), distributed by arXiv under the Creative Commons Attribution 4.0 International licence, http://creativecommons.org/licenses/by/4.0/, as recorded in the arXiv licence metadata for this exact version and shown on the abstract page https://arxiv.org/abs/2602.16126v5. Full licence notice: \texttt{licenses/arxiv-2602.16126-CC-BY-4.0.txt}.\par\smallskip

\noindent\textit{Editorial scope.} Counted unit: the article's proposition prop:Lambda\_heat\_kernel\_mixture (source0.tex lines 304--321). The article's complete original proof of it is delivered with it (source0.tex lines 322--348, one \texttt{proof} environment of 27 lines). No mathematics is written, repaired or re-derived: the statement and the proof are the article's own lines, in the article's own notation, and no retained line is substituted or re-flowed.  Retained uncounted as the source-local context the proof needs: lines 161--163.  Nothing was omitted from any retained span, so the package carries no omission record.  No retained line contains a citation, so the wrapper carries no bibliography and no bibliography entry.  No retained line contains a figure environment, an image-inclusion command or drawing code, so no image is delivered and the package declares no asset dependency.  Editorial formatting only, in addition: an AMS \texttt{amsart} preamble that restates the article's own declarations and macros used by the retained lines, namely the environments \texttt{assumption}, \texttt{proposition} on separate counters, together with the standard packages those lines need.  The retained environments print in this document as Assumption 1, Proposition 1 in order of appearance, whereas the article prints the same items with the numbering of its own full document.  The complete native source of the article as distributed in the arXiv e-print for this exact version is bundled unchanged as \texttt{sources/194849/source0.tex}.
\begin{assumption}\label{assump:WD}
    Let $p_t$ denote the heat kernel, i.e. the kernel of the heat semigroup $P_t$, which we will assume exists. We assume there exists $\Lambda\in(0,\infty)$ such that \[\sup_{x,x'\in M}\int_0^{\infty}\iint_{M^2}p_t(x,y)p_t(x',y')|R(y,y')|dydy'dt=\Lambda.\]
\end{assumption}

\begin{proposition}[Heat-kernel mixture covariances satisfy Assumption~\ref{assump:WD}]\label{prop:Lambda_heat_kernel_mixture}
Let $\nu$ be a finite nonnegative Borel measure on $\RR_+$ and define
\[
R_\nu(y,y') := \int_{\RR_+} p_s(y,y')\,\nu(ds).
\]
Then $R_\nu$ is symmetric and nonnegative definite.

Moreover, if
\begin{equation}\label{eq:Hnu_condition}
\int_0^\infty H(s)\,\nu(ds) < \infty,
\qquad
H(s):=\int_s^\infty \sup_{z\in M} p_u(z,z)\,du,
\end{equation}
then Assumption~\ref{assump:WD} holds for $R_\nu$, and in fact if $\nu\neq 0$
\[
0<\Lambda=\Lambda(R_\nu)\le \frac12 \int_{[0,\infty)} H(s)\,\nu(ds).
\]
\end{proposition}

\begin{proof}
Positivity is immediate, since each $p_s(\cdot,\cdot)$ is positive definite and $R_\nu$ is a nonnegative mixture of such kernels.

For the Assumption~\ref{assump:WD} quantity, by Fubini and the semigroup property,
\begin{align*}
&\int_0^\infty \iint_{M^2} p_t(x,y)\,p_t(x',y')\,R_\nu(y,y')\,dydy'\,dt \\
=&
\int_0^\infty \int_0^\infty \iint_{M^2} p_t(x,y)\,p_t(x',y')\,p_s(y,y')\,dydy'\,dt\,\nu(ds) \\
=&
\int_0^\infty \int_0^\infty p_{2t+s}(x,x')\,dt\,\nu(ds).
\end{align*}
Using $p_u(x,x')\le \sup_{z\in M}p_u(z,z)$ and the change of variables $u=2t+s$ gives
\[
\int_0^\infty p_{2t+s}(x,x')\,dt
\le
\frac12\int_s^\infty \sup_{z\in M} p_u(z,z)\,du
=
\frac12 H(s).
\]
Taking the supremum over $x,x'$ yields
\[
\Lambda(R_\nu)\le \frac12\int_0^\infty H(s)\,\nu(ds),
\]
which is finite by \eqref{eq:Hnu_condition}.

Finally, if $\nu\neq 0$, then $R_\nu\not\equiv 0$, and since the heat kernel is nonnegative, the integrand in Assumption~\ref{assump:WD} is positive on a set of positive measure. Hence $\Lambda(R_\nu)>0$.
\end{proof}

\end{document}
